% P19-E3: the complete reported task families in Tversky-Kahneman 1983.
\chapter{Conjunction Judgments and Event Inclusion}
\label{sec:p19-conjunction-judgments}

A more detailed account may fit a description better while covering fewer
possible outcomes. Probability and descriptive fit therefore need not
rank the same statements alike. To assess a conjunction judgment, first
identify the events, the information held fixed, and the question the
respondent understood.

This chapter presents the task families and numerical comparisons reported
in Tversky and Kahneman's original 1983 article
\cite{TverskyKahneman1983Conjunction}. It follows the evidence-record
distinctions of \cref{def:p19-evidence-record}. The profiles and scenarios
below are independently worded descriptions of the source tasks.
Their historical forecasts, clinical vignettes and fictional allegations
are study materials, not current forecasts, medical guidance or claims
about actual people. The article reports a selected series of demonstrations;
its percentages are not estimates of how often people err in ordinary life.

\section{The probability benchmark and its interpretation}
\label{sec:p19-conjunction-benchmark}

\begin{hbdefinition}{def:p19-conjunction-events}{A common event space and information set}
Let \((\Omega,\mathcal F,P)\) be a probability space and let
\(A,B,E\in\mathcal F\), with \(P(E)>0\).
The conjunction is the intersection \(A\cap B\).
This writes the source's \(A\mathbin{\&}B\) as an event intersection.
The information \(E\) is held fixed when comparing
\(P(A\cap B\mid E)\) and \(P(B\mid E)\).
An elicited numerical judgment \(q(C)\) about event \(C\) is a response;
it is not assumed to satisfy the probability axioms.
A reported conjunction reversal is the strict ordering
\(q(A\cap B)>q(B)\), interpreted on a common increasing scale.

In a likelihood-ranking task, smaller numerical ranks mean greater
judged probability, so the corresponding reversal has
\(\operatorname{rank}(A\cap B)<\operatorname{rank}(B)\).
A rating of similarity, typicality or informativeness is a different
measure, even if the task uses the same list of descriptions.
\index{conjunction fallacy}\index{event inclusion}
\end{hbdefinition}

\begin{hblemma}{thm:p19-conjunction-benchmark}{proved=annex:p19.conjunction-benchmark}{Inclusion, conditioning and frequency}
Under the definition above,
\[
             P(A\cap B\mid E)\leq
                    \min\{P(A\mid E),P(B\mid E)\}.
\]
Equality with \(P(B\mid E)\) holds exactly when
\(P((B\setminus A)\cap E)=0\). No independence assumption is needed.
For a finite collection of cases, the count satisfying both attributes
is at most the count satisfying either attribute. The same inequality
holds for frequencies when the positive denominator is the same.
Upper and lower probability envelopes over one fixed nonempty family
of probability measures are also monotone under event inclusion.
\end{hblemma}

\begin{hbtheorem}{thm:p19-conjunction-interpretations}{proved=annex:p19.conjunction-interpretations}{Three questions that are not interchangeable}
The following comparisons can all hold in a coherent probability model:
\[
 P(A\cap B)>P(A^c\cap B),\qquad
 P(B\mid A)>P(B),\qquad
 P(E\mid A\cap B)>P(E\mid B).
\]
They do not contradict the conjunction benchmark.
For example, assign the four cells
\((A\cap B,A\cap B^c,A^c\cap B,A^c\cap B^c)\)
probabilities \((.20,.05,.10,.65)\). Split each cell into \(E\)
and \(E^c\), using an \(E\)-fraction \(.9\) in each \(A\) cell
and \(.1\) in each \(A^c\) cell.
Then
\[
 \begin{gathered}
 P(B)=.30,\quad P(B\mid A)=.80,\quad P(A\cap B)=.20,\\
 P(E)=.30,\quad P(E\mid A\cap B)=.90,\quad P(E\mid B)=19/30,\\
 P(A\cap B\mid E)=.60\leq P(B\mid E)=19/30.
 \end{gathered}
\]
The same two marginal probabilities do not determine a conjunction:
with \(P(A)=P(B)=1/2\), the intersection may have any probability
between zero and \(1/2\), depending on the joint law.
\end{hbtheorem}

Consequently, interpreting an unqualified \(B\) as \(B\cap A^c\),
answering a conditional question instead of a joint question, or
evaluating the fit of evidence to a hypothesis changes the target.
An empirical analysis must investigate those possibilities rather
than decide the respondent's interpretation from the answer alone.
The finite construction establishes logical possibilities; it does
not establish which interpretation any participant actually used.

\section{What the experimental comparisons establish}
\label{sec:p19-conjunction-designs}

The source distinguishes three procedures
\cite[pp.~297--298]{TverskyKahneman1983Conjunction}.
In an indirect test, different respondents judge the conjunction and
the less representative constituent. In a direct-subtle test, one
respondent encounters both, but their inclusion relation is not
emphasized. A direct-transparent test makes that relation more explicit.
These procedures test different opportunities to recognize inclusion.

The article calls a reversal in a person's direct comparison a
\emph{conjunction fallacy}, and a reversal inferred from separate-group
comparisons a \emph{conjunction error}. The terminology identifies an
operational comparison; it does not diagnose a person's reasoning
in every task.

\begin{hblemma}{thm:p19-conjunction-aggregation}{proved=annex:p19.conjunction-aggregation}{Common respondents and separate-group summaries}
If each of finitely many respondents gives numbers \(q_i\leq p_i\),
their averages with the same weights \(w_i\geq0\),
\(\sum_iw_i=1\), satisfy
\(\sum_iw_iq_i\leq\sum_iw_ip_i\).
A reversed common-weight mean therefore requires at least one
individual reversal. Different respondent groups or different weights
do not support that inference. Two individually coherent respondents
may have \((q_1,p_1)=(.8,.9)\) and
\((q_2,p_2)=(.1,.2)\); reporting the first respondent's conjunction
and the second's constituent gives \(.8>.2\).
\end{hblemma}

Separate-group random sampling can support an estimate of a systematic
elicitation effect, with its sampling uncertainty, but it does not
recover each respondent's unasked comparison. Mean ranks, ordinal
rating scores, percentages choosing an alternative and geometric
means of probability estimates are different summaries. They should
not be pooled as though they were measurements on one cardinal scale.

\section{Availability and representativeness checks}
\label{sec:p19-conjunction-availability}

\paragraph{Word retrieval and estimated frequency.}
UBC students had 60 seconds to produce seven-letter words.
The narrower condition required the suffix \emph{ing}; the broader
condition required an \emph{n} in position six. Every word meeting
the first requirement meets the second. Mean numbers produced were
6.4 and 2.9, with reported \(t(44)=4.70,\ p<.01\).
The article does not give the separate group sizes for that retrieval task.
Production during a timed search is not a census of the available words.

Separate frequency questions described four novel pages, approximately
2,000 words, and used response categories
0, 1--2, 3--4, 5--7, 8--10, 11--15 and 16 or more.
The source-reported median estimates were 13.4 for the suffix condition
(\(N=52\)) and 4.7 for the sixth-position condition (\(N=53\));
the reported median test had \(p<.01\).
A parallel seven-letter comparison used suffix \emph{ly} versus
\emph{l} in position six, with reported medians 8.8 and 4.4.
No group sizes or test statistic are supplied for that parallel comparison
\cite[p.~295]{TverskyKahneman1983Conjunction}.
The smaller class was easier to retrieve under the given search cue;
the numerical frequency comparison was between groups.

\paragraph{Typicality and estimated frequency can dissociate.}
Among Stanford undergraduates, 68\% of a group of 105 selected a
narrow female-student weight band, 124--125 pounds, as more representative
than the broad tail above 135 pounds. In a different group of 102,
78\% selected the broad tail as more frequent.
For the Hollywood-actress items, 65\% of the 105-person group judged
more than four divorces more representative than Democratic voting,
while 83\% of the 102-person group judged Democratic voting more frequent.
These are reported judgments about the source's category descriptions,
not verified population frequencies. Neither pair consists of nested
events, so these differences are not themselves conjunction violations
\cite[p.~296]{TverskyKahneman1983Conjunction}.

\section{Person descriptions: Bill and Linda}
\label{sec:p19-conjunction-persons}

The Bill vignette described an intelligent 34-year-old man, strong in
mathematics and weak in social studies and humanities, with little
imagination, a compulsive disposition and little animation.
It was designed to favor the accountant stereotype over a jazz
hobbyist stereotype. Let \(A\) denote accountant and \(J\) a jazz hobby.
Its critical conjunction was \(A\cap J\).
The Linda vignette described an unmarried 31-year-old woman as outspoken
and highly intelligent, with a philosophy major, concern about
discrimination and social justice, and participation in antinuclear
demonstrations. It was designed to favor feminist activity \(F\)
over bank-teller work \(T\); its critical conjunction was \(T\cap F\).
These intended associations are experimental constructions, not
assertions that professional or political categories have those traits.

Each list contained the two constituents, their conjunction, and five
fillers. Bill's fillers concerned a physician with a poker hobby,
architecture, surfing, reporting and mountaineering. Linda's concerned
elementary teaching, bookstore work with yoga, psychiatric social work,
League of Women Voters membership and insurance work.

\paragraph{Representativeness task.}
A group of 88 UBC undergraduates ranked the eight descriptions by
resemblance to a typical category member. The predicted ordering
\(A>A\cap J>J\) occurred for 87\%; the ordering
\(F>T\cap F>T\) occurred for 85\%.
These are representativeness rankings, not probability-error rates
\cite[p.~297]{TverskyKahneman1983Conjunction}.

\paragraph{Direct and indirect probability rankings.}
The source's naive group comprised Stanford and UBC undergraduates
without probability/statistics training; its informed group comprised
first-year psychology/education graduates and medical students at Stanford
with introductory statistical knowledge; its sophisticated group comprised
Stanford decision-science doctoral students with advanced coursework.
These were educational groups, not randomized training interventions.

Each participant first ranked all eight alternatives for one person
in a direct test, then received the other person in an indirect format
containing either the conjunction or its constituents. Five common
fillers anchored the comparisons. Table~\ref{tab:p19-conjunction-ranks}
retains every numerical row of the source table. The ranks are relative
to the common filler set; smaller numbers indicate greater judged likelihood.
\(C\) denotes the critical conjunction and \(B\) its less representative
constituent. \(N_i\) is the total indirect sample, divided approximately
equally between its two versions, not the sample size of each version.

\begin{table}[htbp]
\centering\small
\setlength{\tabcolsep}{3pt}
\begin{tabular}{llrrrrrrr}
\hline
Group & Person & \(N_d\) & \(V_d\%\) & \(R_d(C)\) & \(R_d(B)\)
      & \(N_i\) & \(R_i(C)\) & \(R_i(B)\)\\
\hline
Naive & Bill & 94 & 92 & 2.5 & 4.5 & 88 & 2.3 & 4.5\\
Naive & Linda & 88 & 89 & 3.3 & 4.4 & 86 & 3.3 & 4.4\\
Informed & Bill & 56 & 86 & 2.6 & 4.5 & 56 & 2.4 & 4.2\\
Informed & Linda & 53 & 90 & 3.0 & 4.3 & 55 & 2.9 & 3.9\\
Sophisticated & Bill & 32 & 83 & 2.6 & 4.7 & 32 & 2.5 & 4.6\\
Sophisticated & Linda & 32 & 85 & 3.2 & 4.3 & 32 & 3.1 & 4.3\\
\hline
\end{tabular}
\caption{Source Table 1: direct reversal rates and common-filler mean ranks.
The two formats include repeated contributions by participants.}
\label{tab:p19-conjunction-ranks}
\end{table}

The conjunction received the higher likelihood position in all twelve
direct/indirect row comparisons. The source reports 88\% overall direct
reversals and correlations between aggregate probability and
representativeness ranks of .96 for Bill and .98 for Linda.
These correlations are between mean rank profiles, not individual-level
process measurements. The row sample sizes must not be added to claim
that the direct and indirect stages used entirely separate people
\cite[p.~298]{TverskyKahneman1983Conjunction}.

\section{Transparent Linda variants}
\label{sec:p19-conjunction-transparent}

The following tasks used separate groups and different response measures
\cite[pp.~299--300]{TverskyKahneman1983Conjunction}.

\begin{itemize}
\item \textbf{Two-option probability choice.}
      Among 142 UBC undergraduates, 85\% chose \(T\cap F\) over \(T\).
      Reversing the alternative order for approximately half the group
      had no reported detectable effect.
\item \textbf{Nine-point ratings.}
      Among 119 UBC undergraduates, mean ratings were 3.5 for \(T\)
      and 5.6 for \(T\cap F\); 82\% rated the conjunction higher.
      The endpoints of the scale were extremely unlikely and extremely
      likely. These ratings are not probabilities .35 and .56.
\item \textbf{Choosing an argument.}
      Of 58 undergraduates, 65\% preferred an argument appealing to
      resemblance over the inclusion argument. This measures selection
      between reasons, not a numerical probability estimate.
\item \textbf{Explicit inclusive wording.}
      The teller statement was expanded to include both feminist and
      nonfeminist tellers. For 75 respondents, mean likelihood ratings
      were 3.8 for this inclusive version \(T^*\) and 5.1 for \(T\cap F\),
      with reported \(p<.01\) by a \(t\)-test. The conjunction was rated
      higher by 57\% and lower by 16\%.
\item \textbf{Separate-group wording comparison.}
      After a common psychiatric-social-worker filler, mean ratings
      were 3.3 for \(T\), 3.9 for \(T^*\), and 4.5 for \(T\cap F\).
      The paper reports each mean as significantly different from both
      others, but supplies no group sizes or exact test values here.
\item \textbf{Hypothetical betting choice.}
      With the same \$10 reward attached to either selected event,
      56\% of 60 undergraduates chose the narrower event.
\item \textbf{Graduate rating group.}
      Among 64 social-science graduates at Berkeley and Stanford with
      several statistics courses, means were 3.8 for \(T\) and 3.5
      for \(T\cap F\), and 36\% gave a conjunction reversal.
      This is a different sample and format from the advanced
      decision-science rows in Table~\ref{tab:p19-conjunction-ranks}.
\end{itemize}

These results make a single, context-free claim about the effect of
statistical education inappropriate. The main ranking groups showed
similar reversal patterns, whereas the graduate transparent-rating
group showed a different pattern. The group comparisons do not by
themselves estimate the causal effect of taking a statistics course.

\section{Clinical and sporting judgments}
\label{sec:p19-conjunction-expert-sport}

\paragraph{Clinical predictions.}
Two main groups comprised 37 Boston-area internists in a Harvard
postgraduate course and 66 internists with New England Medical Center
admitting privileges. They judged five constructed clinical problems.
Each participant received two or three indirect problems first and
the remaining problems directly; item allocation balanced the formats.
The wording asked which conditions would be among those experienced
and explicitly permitted several conditions.

The one clinical vignette printed in full concerns a 55-year-old woman
with angiographically confirmed pulmonary embolism ten days after
gallbladder removal. Its critical comparison
adds dyspnea to hemiparesis; the consultants had selected the former
as representative and the latter as unrepresentative for this task.
The other listed alternatives included calf pain, pleuritic pain,
hemoptysis and the syncope/tachycardia combination. The article does
not print the full stimuli for the other four problems.

An independent Stanford group of 32 physicians ranked representativeness.
Across all five problems, correlations of aggregate probability and
representativeness ranks exceeded .95. Overall mean ranks for conjunction
and less likely constituent were 2.7 and 4.6 in direct tests, and
2.8 and 4.3 indirectly. Direct reversal rates ranged from 73\% to 100\%,
averaging 91\%.
In the end-of-questionnaire interpretation check, 60 of 62 physicians
endorsed an inclusive, multiple-symptom reading. A further group of
24 physicians, mostly Stanford residents, discussed their responses;
the report describes surprise and lack of defense but gives no coded
interview distribution. The 62-person check is not an additional
independent trial series \cite[pp.~301--302]{TverskyKahneman1983Conjunction}.

\paragraph{A tennis forecast and a separate interpretation check.}
In October 1980, 93 people recruited through a University of Oregon
newspaper advertisement forecast a hypothetical Wimbledon final
conditional on Borg reaching it in 1981. Mean ranks were 1.7 for
winning the match, 2.7 for losing the first set, 2.2 for losing that
set but winning the match, and 3.5 for winning the first set but
losing the match. The third event was ranked above the second by 72\%.

A separate group of 56 statistically naive Stanford respondents classified
win/loss sequences against these descriptions. The report says that,
with few exceptions, their classifications respected the intended
inclusion. For example, a first-set loss remained an instance of that
event whether the player later won or lost the match.
This checks interpretation in that group; it does not directly measure
the interpretation of each of the 93 forecasters
\cite[pp.~302--303]{TverskyKahneman1983Conjunction}.

\section{A twenty-roll betting task}
\label{sec:p19-conjunction-betting}

The task used a six-faced die with four green faces and two red faces,
rolled twenty times. A chosen consecutive string paid \$25 if it
appeared anywhere in those rolls. The three strings were
\[
       s_1=\texttt{RGRRR},\qquad
       s_2=\texttt{GRGRRR},\qquad
       s_3=\texttt{GRRRRR}.
\]
The critical inclusion is that every occurrence of \(s_2\) contains
an occurrence of \(s_1\); it does not require evaluating their
representativeness. The third string is a comparison option, not
another event shown to be a subset of the first by this argument.

\begin{hbtheorem}{thm:p19-conjunction-bets}{proved=annex:p19.conjunction-bets}{Nested payouts and the actual twenty-roll events}
Let \(E_s\subseteq\{R,G\}^{20}\) mean that string \(s\) occurs
consecutively somewhere in the twenty rolls. Then
\[
 E_{\texttt{GRGRRR}}\subseteq E_{\texttt{RGRRR}},
 \qquad
 E_{\texttt{RGRRRG}}\subseteq E_{\texttt{RGRRR}}.
\]
At the same nonnegative prize and same cost, the shorter-string
payout is at least the longer-string payout in every outcome.
With independent rolls, \(P(G)=2/3,\ P(R)=1/3\), and a strictly
positive prize, the expected-money inequalities are strict.

The exact occurrence probabilities in this model are
\[
\begin{array}{c|r|c}
s&3^{20}P(E_s)&P(E_s)\text{ approximately}\\ \hline
\texttt{RGRRR}&435643544&.12494135\\
\texttt{GRGRRR}&282259980&.08095137\\
\texttt{GRRRRR}&142308810&.04081377\\
\texttt{RGRRRG}&277222020&.07950650
\end{array}
\qquad 3^{20}=3486784401.
\]
The proof counts overlapping possible occurrences; multiplying the
probabilities of the letters gives only a specified-position probability.
These exact numerical values are handbook calculations, not participant
responses reported in the source.

Unequal rewards change the decision question: for nested events
with probabilities .4 and .6, rewards two and one give expected
payoffs .8 and .6. Preferring the narrower event with that different
reward is compatible with correct probability ordering.
\end{hbtheorem}

The reported results were as follows
\cite[pp.~303--304]{TverskyKahneman1983Conjunction}.
A representativeness group of 50 placed \(s_2\) first and \(s_3\)
last in 88\% of responses.
In the choice study, 260 UBC and Stanford students participated in
classroom groups of 30--50. Of these, 125 faced real payoffs and
the rest hypothetical payoffs. The dominated \(s_2\) was selected
by 65\% and 62\%, respectively; \(s_3\) was selected by 2\% in
each condition. The report pooled the universities after finding no
significant difference, without reporting a full effect estimate.

A new hypothetical-choice group of 59 saw
\(\texttt{RGRRRG}\) in place of \(s_2\), so the common prefix
made inclusion more visible; 63\% nevertheless chose it over \(s_1\).
Another group of 88 judged two supplied arguments, rather than
choosing a gamble: 76\% selected the valid inclusion/dominance reason
over the reason based on matching the die's green/red proportions.
The incentive, alternative-order and argument-choice variants are
different experiments; their percentages are not interchangeable
measures of one underlying error rate.

\section{Causal additions and coherent stories}
\label{sec:p19-conjunction-causes}

The source contrasts two task structures. A profile \(M\) can make one
added property representative while making the basic outcome
unrepresentative. Alternatively, an added event can provide a
plausible cause of the basic outcome. Both diagrams in its Figure 1
give a negative association from \(M\) to the basic event \(B\).
The first adds a positive association from \(M\) to \(A\); the second
adds one from \(A\) to \(B\). These are association diagrams;
they do not supply structural causal models with intervention semantics.

\paragraph{Health-survey probability choices.}
A hypothetical survey sample covered adult men in British Columbia
across ages and occupations. Among 115 Stanford and UBC undergraduates,
58\% judged a randomly selected man's history of heart attack combined
with age over 55 more probable than the heart-attack history alone.
In a version assigning the age condition to a second, unrelated
randomly selected man, 29\% of 90 respondents reversed the ordering.
The version changes who bears the second property, not the
mathematical inclusion of the joint event in the first.
Its intended independence model is a stipulation of the comparison,
not a demonstrated property of a real health survey
\cite[pp.~305--306]{TverskyKahneman1983Conjunction}.

\paragraph{Running times.}
For a fictional college runner with a prior best mile of 4:06,
96 Stanford undergraduates ranked outcomes including a mile below
4:06, a mile below 4:00, a second half-mile below 1:55, their
conjunction, and a first half-mile below 2:05.
The conjunction was ranked above at least one constituent by 76\%
and above both by 48\%. No respondent ranked the sub-four-minute
mile above the broader sub-4:06 outcome.
The different comparisons within the same task therefore had
different outcomes \cite[p.~306]{TverskyKahneman1983Conjunction}.

\paragraph{A motive added to an alleged crime.}
In a fictional investigation vignette, 171 UBC students ranked
allegations about a 42-year-old married businessman with two children.
The fictional profile describes him as mild and somewhat secretive,
operating a New York import-export business and travelling frequently
to Europe and the Far East. A prior conviction concerned smuggling
precious stones and metals, including uranium, with a six-month
suspended prison sentence and a large fine.
One group of 86 received an employee-homicide item, while a group
of 85 received the item with an added witness-silencing motive.
Other allegations concerning sexual offending, espionage/secrets and
drug addiction served as common comparisons.
The motive-qualified item had mean rank 2.90 versus 3.17 for the
unqualified item, with reported \(p<.05\) by a \(t\)-test.
It was ranked above the drug-addiction comparison by 50\%, compared
with 23\% for the unqualified item. This is a between-group error
comparison, not a recorded within-person conjunction reversal.
A second police-officer example is described qualitatively without
sample sizes or numerical results \cite[pp.~306--307]{TverskyKahneman1983Conjunction}.

\paragraph{Historical forecasting scenarios.}
In April 1982, 245 UBC undergraduates estimated events specified
for 1983. Separate versions asked about a North American flood
with more than 1,000 drowning deaths or a California earthquake causing
a flood with that death count. The paper treats the latter as a
scenario for the same broad flood event; that comparison presupposes
the same geographic and outcome interpretation.
Their source-reported geometric mean estimates were
2.2\% for the flood and 3.1\% for the earthquake/flood scenario,
with reported \(p<.01\) by a Mann--Whitney test.

At the July 1982 International Congress on Forecasting in Istanbul,
115 participants, mostly professionally involved in forecasting and
planning, received parallel scenario tasks. Geometric mean estimates
were .14\% for suspension of US--Soviet diplomatic relations and
.47\% for a Soviet invasion of Poland followed by that suspension
(\(p<.01\), Mann--Whitney). For a 30\% decline in US oil consumption,
the estimate was .22\%, versus .36\% when a dramatic price increase
was added; the paper calls this difference significant without
providing the corresponding test value.

Both studies used nine probability categories, from below .01\%
through .1, .5, 1, 2, 5, 10 and 25\%, to 50\% or more.
The paper gives study totals but not the exact per-version sample
sizes for each item. These are judgments elicited in 1982,
not present-day forecasts, and the reported comparisons do not
score them against later outcomes
\cite[pp.~307--308]{TverskyKahneman1983Conjunction}.

\section{Frequency wording and additional category estimates}
\label{sec:p19-conjunction-frequencies}

The health-survey item was also studied with numerical estimates
\cite[pp.~308--309]{TverskyKahneman1983Conjunction}.
Here the same respondent estimated the broad heart-attack category
and the joint heart-attack/over-55 category.

\begin{table}[htbp]
\centering\small
\begin{tabular}{p{.65\linewidth}rr}
\hline
Task version & \(N\) & Reversals\\ \hline
Percentage estimates, statistically naive UBC sample &147&65\%\\
Percentage estimates, advanced Stanford undergraduates with statistics &62&28\%\\
Percentage estimates, with the over-55 percentage requested first &159&31\%\\
Counts among 100 hypothetical survey participants &117&25\%\\
Counts among 100, with the over-55 count requested first &360&11\%\\ \hline
\end{tabular}
\caption{Reported within-person ordering reversals in five survey-task versions.
Different rows represent different samples or procedures.}
\label{tab:p19-conjunction-frequency}
\end{table}

For the first row, mean estimates were 18\% for heart-attack history
and 30\% for the conjunction. Reversing the order of the two
constituents within the conjunction had no reported significant effect.
A separate group of 55 answered the explicitly conditional question
about heart-attack history among men over 55, with mean estimate
59\%. The original error-making subgroup's mean conjunction estimate
was 45\%, not 59\%. That selected subgroup is part of the original
sample, not another independent cohort.

Counts and percentages must be compared with a common denominator.
The 100 in the task is the hypothetical survey size, whereas the
\(N\)'s in Table~\ref{tab:p19-conjunction-frequency} are numbers
of respondents. The table shows that formulation and auxiliary
questions mattered in these samples. It does not identify a unique
psychological mechanism for the reduction, establish that frequency
wording always eliminates reversals, or supply a population-wide
causal effect of statistical education.

\section{Interpretation, limits and practical use}
\label{sec:p19-conjunction-interpretation}

The article also reports that omitting Linda's personality sketch
and retaining only her age/sex description made almost all respondents
respect inclusion. In the Bill variant, pairing boredom with music
and a jazz hobby reduced conjunction judgments below the constituents.
Its qualitative mathematics/physics example contrasts matching
upper-half or lower-half class positions with mismatched positions
across the two subjects. These examples concern relations between
attributes, which their marginal probabilities do not determine.
These observations lack numerical sample details in the reported passages.
They belong to the source's qualitative interpretation, not an
additional reconstructed dataset.

The authors also distinguish a conceptual fallacy from miscommunication
through respondents' later willingness to acknowledge a substantive,
potentially repeatable mistake, a conceptual rather than merely verbal
error, and an available rule or procedure they should have used.
This is their proposed interpretive criterion, not a validated diagnostic
scale or a substitute for recording the actual response.

The authors consider resemblance, availability and causal coherence;
deliberate substitution, anchoring and failure to distinguish the
requested judgment; and alternatives involving exclusive readings,
conditional or inverse probabilities, informational value and surprise.
The experiments do not identify one mechanism from a response
percentage alone. A communicatively useful detailed message may have
greater value despite lower probability, just as different rewards can
change the preferred bet. Those are different questions from the
equal-payoff probability comparison.

The paper reports that its broader project involved more than 3,000
participants and many problem variants, but explicitly presents only
a selected, nonexhaustive set. The tasks were constructed to elicit
the effect. Missing per-condition sample sizes, raw responses,
uncertainty intervals and unprinted stimuli must not be supplied by
guessing. Rounded percentages do not determine exact participant
counts, and mean-rank correlations do not establish individual
processing strategies. The statistical tests above are the source's
reported analyses, not new reanalyses.

The 1999 language and response-format experiments are supplied in
\cref{sec:p19-conjunction-language}. The wider literature includes separate semantic,
pragmatic, frequency-format and process-model investigations.
This historical presentation does not replace those later source
treatments or claim a present consensus from one article.
Coherence of a set of judgments also differs from calibration against
outcomes, as defined in \cref{sec:p19-calibration-extreme-confidence}.

\begin{hbmethod}{def:p19-conjunction-method}{Analyzing a conjunction judgment}
Write the target events and information explicitly. Check that one
event really includes the other under the intended reading, with
the same person, time horizon, reference class and denominator.
For frequency questions, identify the population being counted.
For a betting question, specify the complete payout conditions.

Record whether the response concerns probability, frequency,
typicality, explanatory fit, surprise or a preferred action.
Use an interpretation check where the language may permit a
different event partition, while recognizing that the check can
itself change the task. Separate individual comparisons from
comparisons of different groups.

If the intended nested-event reading is established, enforce its
probability ordering before using the estimates in a larger inference.
The ordering alone does not choose the correct numerical probabilities
or a unique repair to inconsistent estimates. Retain the measured
response and its provenance when studying behavior; use an explicitly
justified coherent model when performing a probability calculation.
\end{hbmethod}
